% =====================================================================================================
% appB_wilson.tex —— 附录 B  Wilson 线：Feynman 规则、程函恒等式与尖点（writer W05）
% 源：FINAL_OUTLINE_v2 App B；research/c_wilson（DERIVATION.md §2--§3）；VF-WL-1/3/4/13，VF-CONV-23，VF-SCET-16；
%     验证脚本 writing/work/W05/s4_wilson.py、s5_cusp.py（日志 logs/）。
% =====================================================================================================
\chapter{Wilson 线：Feynman 规则、程函恒等式与尖点}\label{app:appB}

本附录证明第~3~章的\cref{prop:prelim-wilson-covariance}、\cref{prop:prelim-wilson-rules}（即\cref{prop:appB-rules}）与\cref{lem:prelim-eikonal-identity}（即\cref{lem:appB-shuffle}），
并给出半无限 Wilson 线的 Ward 恒等式（附录~F）、非阿贝尔指数化（\cref{prop:appB-webs}）与单圈 cusp 反常维度（\cref{prop:appB-cusp-oneloop}）。

% -----------------------------------------------------------------------------------------------------
\section{定义与约定}\label{sec:appB-definition}

沿路径 $z(s)$（$s\in[0,1]$，$z(0)=a$，$z(1)=b$）的 Wilson 线（\cref{def:prelim-wilson}）为
\begin{equation}\label{eq:appB-W}
W[b,a]=P\exp\Big(\ii g\int_a^b\dd z\cdot A\Big)=\sum_{m=0}^\infty(\ii g)^m\!\!\int\limits_{0<s_1<\dots<s_m<1}\!\!\dd s_1\cdots\dd s_m\;
\dot z\cdot A(z(s_m))\cdots\dot z\cdot A(z(s_1)),
\end{equation}
$A\equiv A^at^a$，离 $b$ 近的点在左；$W(s)\equiv W[z(s),a]$ 是 $\frac{\dd}{\dd s}W=\ii g\,\dot z\cdot A(z(s))\,W$、$W(0)=1$ 的唯一解。

\emph{\cref{prop:prelim-wilson-covariance}的证明。}由于 $W[c,b]W[b,a]$ 与 $W[c,a]$ 满足同一方程与同一初值，二者相等；取 $c=a$ 得 $W[a,b]W[b,a]=1$，对厄米的 $A$ 有 $W[b,a]^\dagger=W[a,b]$。
在 $D=\partial-\ii gA$ 的规范变换 $A'=UAU^{-1}-\frac{\ii}g(\partial U)U^{-1}$ 下，$V(s)=U(z(s))W(s)U^{-1}(a)$ 满足
\[
\frac{\dd V}{\dd s}=\Big[\frac{\dd U}{\dd s}U^{-1}+\ii g\,U\dot z\cdot AU^{-1}\Big]V=\ii g\,\dot z\cdot A'\,V,\qquad V(0)=1,
\]
故 $W[b,a]\to U(b)W[b,a]U^{-1}(a)$（非齐次项反号则不成立；SU(3) 中数值核对）。% check: s4_wilson.py (d)
伴随线 $W^{(A)}_{ab}=2\Tr[t^aW[b,a]t^bW[a,b]]$ 即 $W t^bW^{-1}=t^aW^{(A)}_{ab}$。

\emph{约定。}(C1) $A(x)=\int\frac{\dd^4k}{(2\pi)^4}\ee^{-\ii k\cdot x}\tilde A(k)$，$k$ 流入 Wilson 线；(C2) $n^0>0$，$W[\infty n,y]=\lim_{L\to\infty}W[y+Ln,y]$ 等，端点 $y\ne0$ 时乘 $\ee^{-\ii y\cdot K_{\rm tot}}$；
(C3) 绝热正规子：半无限线上的场乘以 $\ee^{-\epsilon s}$（$s$ 为到有限端点的仿射距离），最后 $\epsilon\to0^+$。Collins 用 $D=\partial+\ii g_0tA$\cite[式~(2.2)]{Collins2011}，故 $g_{\rm Collins}=-g$，
其 Wilson 线 $P\exp(-\ii g_0\int n\cdot A)$\cite[式~(13.40)]{Collins2011}作为算符与我们的相同；他的图~7.10 与式~(11.4)--(11.6) 须按作者勘误使用（附录~Q 第~19~条）。

% -----------------------------------------------------------------------------------------------------
\section{动量空间 Feynman 规则}\label{sec:appB-rules}

本节给出第~3~章\cref{prop:prelim-wilson-rules}的证明。

\begin{proposition}[Wilson 线的 Feynman 规则 \tagR]\label{prop:appB-rules}
（$R_{\mathrm{def}}$：由算符定义直接推出。）按 (C1)--(C3)，直半无限 Wilson 线中连接 $m$ 个胶子（动量 $k_j$ 流入、指标 $\mu_j$、色 $a_j$；记 $n\cdot\tilde A_j\equiv n^{\mu_j}\tilde A^{a_j}_{\mu_j}(k_j)t^{a_j}$，按所写次序相乘）的项为
\begin{align}
W^{(m)}[0,-\infty n]&=\int\!\prod_j\frac{\dd^4k_j}{(2\pi)^4}\,(\ii g)^m\,n\cdot\tilde A_m\cdots n\cdot\tilde A_1\prod_{j=1}^m\frac{\ii}{n\cdot(k_1+\dots+k_j)+\ii j\epsilon},\label{eq:appB-past}\\
W^{(m)}[\infty n,0]&=\int\!\prod_j\frac{\dd^4k_j}{(2\pi)^4}\,(\ii g)^m\,n\cdot\tilde A_m\cdots n\cdot\tilde A_1\prod_{j=1}^m\frac{-\ii}{n\cdot(k_j+\dots+k_m)-\ii(m-j+1)\epsilon},\label{eq:appB-future}\\
W^{(m)}[-\infty n,0]&=\int\!\prod_j\frac{\dd^4k_j}{(2\pi)^4}\,(-\ii g)^m\,n\cdot\tilde A_m\cdots n\cdot\tilde A_1\prod_{j=1}^m\frac{\ii}{n\cdot(k_j+\dots+k_m)+\ii(m-j+1)\epsilon},\label{eq:appB-pastc}\\
W^{(m)}[0,\infty n]&=\int\!\prod_j\frac{\dd^4k_j}{(2\pi)^4}\,(-\ii g)^m\,n\cdot\tilde A_m\cdots n\cdot\tilde A_1\prod_{j=1}^m\frac{-\ii}{n\cdot(k_1+\dots+k_j)-\ii j\epsilon},\label{eq:appB-futc}
\end{align}
% VF-WL-1
其中胶子按路径次序编号：式~\eqref{eq:appB-past} 中 $k_1$ 最靠近 $-\infty$，式~\eqref{eq:appB-future} 与 \eqref{eq:appB-pastc} 中 $k_1$ 最靠近 $0$，式~\eqref{eq:appB-futc} 中 $k_1$ 最靠近 $+\infty$。
等价地（普适规则）：每个连接点一个顶点 $\ii g\,t^av^\mu$（$v$ 为沿取向的切矢量），除伸向无穷远的一段外每段一个 $\ii/(v\cdot\ell+\ii0)$（$\ell$ 沿取向流动，伸向无穷远的一段 $\ell=0$），
色矩阵按路径次序，有限 $\epsilon$ 下某段的 $\ii0$ 为 $\ii\epsilon$ 乘以该段与无穷远端之间的连接点数。单胶子项与“未来减过去”为
\begin{equation}\label{eq:appB-onegluon}
W[\infty n,0]\ni\frac{g\,t^an\cdot\tilde A^a(k)}{n\cdot k-\ii0},\qquad W[-\infty n,0]\ni\frac{g\,t^an\cdot\tilde A^a(k)}{n\cdot k+\ii0},\qquad
\text{差}=2\pi\ii\,g\,t^an\cdot\tilde A^a(k)\,\delta(n\cdot k).
\end{equation}
\uses{def:prelim-wilson}
\end{proposition}

\begin{proof}
以 $W[0,-\infty n]=P\exp(\ii g\int_{-\infty}^0\dd\lambda\,n\cdot A(\lambda n))$ 为例：代入 (C1) 并乘绝热因子，$m$ 阶项的参数积分为
$\int_{-\infty<\lambda_1<\dots<\lambda_m<0}\prod_j\dd\lambda_j\,\ee^{-\ii\lambda_j(n\cdot k_j+\ii\epsilon)}$。先积 $\lambda_1$ 得 $\frac{\ii}{n\cdot k_1+\ii\epsilon}\ee^{-\ii\lambda_2(n\cdot k_1+\ii\epsilon)}$，
它与下一个相因子合并为 $\ee^{-\ii\lambda_2(n\cdot(k_1+k_2)+2\ii\epsilon)}$；归纳地第 $j$ 次积分给出 $\ii/[n\cdot(k_1+\dots+k_j)+\ii j\epsilon]$，$\lambda_m$ 积到 $0$。色矩阵的次序即式~\eqref{eq:appB-W} 的路径次序。
$W[\infty n,0]$ 的区域为 $0<\lambda_1<\dots<\lambda_m<\infty$、阻尼 $\ee^{-\epsilon\lambda_j}$，先积 $\lambda_m$，得式~\eqref{eq:appB-future}；$W[-\infty n,0]$、$W[0,\infty n]$ 的取向为 $-n$，
每个顶点多一个负号，动量仍从无穷远端累加。普适规则是这四式的改写（例如 $W[\infty n,0]$ 第 $j$ 个胶子前一段沿取向的动量为 $-(k_j+\dots+k_m)$）。
式~\eqref{eq:appB-onegluon} 是 $m=1$ 的情形，差由 $\frac1{x-\ii0}-\frac1{x+\ii0}=2\pi\ii\delta(x)$ 给出（即 $\ii g\int_{-\infty}^\infty\dd\lambda\,n\cdot A(\lambda n)$ 的 Fourier 变换）。
四种线的 $m\le3$ 项经符号计算验证；把式~\eqref{eq:appB-W} 在 SU(3) 背景场中数值求解，$g^1$ 至 $g^4$ 的系数与规则相符（残差按 $g^5$ 标度）。% check: s4_wilson.py (a)(b)
\end{proof}
